\documentclass[11pt]{article}
\input{style-sujet}
\usepackage{array}

\begin{document}

\entetesujet{Sujet bac 2016 -- Série D}

% ====================== EXERCICE 1 ======================
\exo{1}{5 points}

Dans le plan complexe $\Cb$ muni d'un repère orthonormal $(O,\vec{u},\vec{v})$, on considère l'application $S$ définie par :
\[ z' = (1 + i)z \]

\begin{enumerate}
  \item Déterminer la nature, le rapport et l'angle de l'application $S$.
  \item Soit le point $A$ d'affixe $z_A = 2i$. Déterminer les affixes des points $B$ et $C$ définis par $S(A) = B$ et $S(B) = C$.
  \item Placer les points $A$, $B$, et $C$ dans un repère du plan.
  \item Soit le point $I$ milieu du segment $[OC]$. Montrer que le triangle $ABI$ est rectangle et isocèle en $B$.
  \item Écrire une équation de troisième degré dont les affixes $z_A$, $z_B$ et $z_C$ définies ci-dessus sont solutions.
\end{enumerate}

% ====================== EXERCICE 2 ======================
\exo{2}{5 points}

L'espace vectoriel $\R^2$ est rapporté à sa base canonique $\mathcal{B} = (\vi,\vj)$. Soit $f$ l'endomorphisme de $\R^2$ défini dans la base $(\vi,\vj)$ par :
\[ f(\vj) = 3\vi - 2\vj \quad \text{et} \quad f \circ f(\vj) = \vj \]

\begin{enumerate}
  \item Calculer $f(\vi)$ et $f \circ f(\vi)$.
  \item En déduire la nature de l'application $f$.
  \item \begin{enumerate}[label=\alph*.]
    \item Qu'est-ce qu'un automorphisme ?
    \item Prouver que $f$ est un automorphisme involutif.
    \item Caractériser l'application $f$.
  \end{enumerate}
  \item Soit les vecteurs $\vec{u} = 2\vi - \vj$ et $\vec{v} = -3\vi + \vj$.
  \begin{enumerate}[label=\alph*.]
    \item Montrer que $\mathcal{B}' = (\vec{u},\vec{v})$ est une base de $\R^2$.
    \item Écrire la matrice de $f$ dans la base $\mathcal{B}'$.
  \end{enumerate}
\end{enumerate}

% ====================== EXERCICE 3 ======================
\exo{3}{7 points}

Dans le plan rapporté à un repère orthonormal $(O,\vi,\vj)$ d'unité graphique 1 cm, on considère la fonction $f$ de la variable réelle $x$ définie sur $\R_+^*$ par $f(x) = \dfrac{e^x}{e^x - 1}$. On note $(C)$, la courbe de $f$ dans le plan.

\begin{enumerate}
  \item Calculer les limites de $f$ en 0 à droite et en $+\infty$.
  \item Déduire que la fonction $f$ admet deux asymptotes que l'on précisera.
  \item \begin{enumerate}[label=\alph*.]
    \item Montrer que pour tout $x$ appartenant à $\R_+^*$, on a $f'(x) = -\dfrac{e^x}{(e^x - 1)^2}$.
    \item Donner le sens de variation de $f$.
    \item Dresser son tableau de variation.
  \end{enumerate}
  \item Tracer la courbe $(C)$ ainsi que ses asymptotes.
  \item Soit $g$ la fonction définie sur $\R_+^*$ par $g(x) = -f(x)$. Construire la courbe $(C')$ de $g$ dans le même repère que $(C)$.
  \item Calculer l'aire $\mathcal{A}$ de la portion du plan délimitée par les courbes $(C)$ et $(C')$, et les droites d'équations $x = 1$ et $x = 2$.
\end{enumerate}

% ====================== EXERCICE 4 ======================
\exo{4}{3 points}

Soit le tableau statistique à double entrée :

\begin{center}
\renewcommand{\arraystretch}{1.3}
\begin{tabular}{|c|c|c|c|}
\hline
$X\,\backslash\,Y$ & 0 & 1 & 2 \\
\hline
$-1$ & 1 & $m$ & 1 \\
\hline
$1$ & 1 & 0 & 2 \\
\hline
$2$ & 2 & 3 & $n$ \\
\hline
\end{tabular}
\end{center}

\begin{enumerate}
  \item Déterminer les lois marginales de $X$ et de $Y$ en fonction de $n$ et $m$.
  \item Déterminer $m$ et $n$ sachant que le point moyen du nuage statistique est $G(1\,;1)$.
  \item On pose $m = 1$ et $n = 1$.
  \begin{enumerate}[label=\alph*.]
    \item Déterminer l'équation de la droite de régression linéaire de $Y$ en $X$ sachant que la covariance de $X$ et $Y$ est égale à $-\dfrac{1}{12}$, la variance de $X$ est $\dfrac{1}{2}$ et celle de $Y$ est $\dfrac{1}{6}$.
    \item Calculer le coefficient de corrélation linéaire entre $X$ et $Y$.
  \end{enumerate}
\end{enumerate}

\end{document}
